% =====================================================================
%  付録：語の対応，記号の一覧，実験の番号と内部の記号，数値の出所，参考文献
%  語の元：paper/sections/*.tex（英語）と paper/ja/sections/*.tex（日本語），slides/CLAUDE.md の「語」の節
%  実験の番号：docs/results.md 3.6 節の「名前」の列．内部の記号：paper/scripts/ledger/comparison.py と paper/CLAIM.md
%  数値の出所の表：figures/appendix_sources.tex（tools/key_sources.py が資料の \NUM{key} を集め，台帳の source 欄から作る）
% =====================================================================
\setband{付録}

\begin{frame}{付録 1：論文の英語の語と，この資料の語（1／4：設定）}
  \rl{語の対応}\label{a:terms}
  \vspace{-0.2em}
  \centering\scriptsize
  \renewcommand{\arraystretch}{1.12}
  \begin{tabularx}{\textwidth}{@{}l l X@{}}
    \toprule
    論文の語 & この資料の語 & 意味 \\
    \midrule
    target / source corpus & ターゲット／ソース & 新しいコーパス／ラベル付きのほかのコーパス \\
    pool & プール & ターゲットのテスト以外の発話．$K$ 本を抜いた残りはラベル無しで使う \\
    $K$ labels per emotion & 感情あたり $K$ 本 & $K=5$ なら 4 感情 × 5 ＝ 20 本 \\
    target-only / source-only / pooled & ターゲットのみ／ソースのみ／合併 & 3 つの学習法 \\
    gain, gain map & 利得，利得の総当たり実験 & 合併 − ターゲットのみ（同じ反復）．すべての順序つきの組で測る \\
    configuration & 種類（世界） & 特徴の組み合わせ A，V，T，AV，VT，AT，AVT \\
    separation $d$ & 識別力 $d$ & 2 つの感情の山の中心の離れ方（山の幅で割る） \\
    threshold offset $o$ & 判断基準のずれ $o$ & 境界が 2 つの中心の真ん中からずれた量 \\
    \bottomrule
  \end{tabularx}
\end{frame}

\begin{frame}{付録 1：論文の英語の語と，この資料の語（2／4：第 I 部）}
  \rl{語の対応}
  \vspace{-0.2em}
  \centering\scriptsize
  \renewcommand{\arraystretch}{1.12}
  \begin{tabularx}{\textwidth}{@{}l l X@{}}
    \toprule
    論文の語 & この資料の語 & 意味 \\
    \midrule
    threshold loss / direction loss & 判断基準による損失／向きによる損失 & 境界を真ん中に動かすと取り戻せる損失／残る損失 \\
    threshold share & 判断基準の割合 & 合併がソースのみを上回った分のうち，$o$ の変化による分の割合 \\
    threshold cost / direction cost & 2 つのコスト & $K$ 本で推定した誤差 $1/(2K)$ と $q/K$ \\
    effective dimension $q$ & $q$ & 向きの誤差の減り方を決める数（音声で約 16〜22） \\
    centering & 中心化 & コーパスごとにラベル無しの発話の平均を引く \\
    hypothesis transfer & 仮説転移 & ソースの分類器から出発して $K$ 本で学び直す \\
    value of a source $n$ & ソースの価値 $n$ & ソースの向きが何本ぶんのターゲットのラベルに当たるか \\
    \bottomrule
  \end{tabularx}
\end{frame}

\begin{frame}{付録 1：論文の英語の語と，この資料の語（3／4：第 II 部）}
  \rl{語の対応}
  \vspace{-0.2em}
  \centering\scriptsize
  \renewcommand{\arraystretch}{1.12}
  \begin{tabularx}{\textwidth}{@{}l l X@{}}
    \toprule
    論文の語 & この資料の語 & 意味 \\
    \midrule
    start, five-value start & 出発点，5 つの値を選んだ出発点 & 全ソースから作る線形分類器 $(W_0, b_0)$．5 つの値は仮のターゲットで選ぶ \\
    anchor & 出発点への縛り & 重みと切片を出発点に近づける L2 罰則 \\
    pseudo-target & 仮のターゲット & ターゲットの代わりにするソースの群（設定の値を選ぶため） \\
    oracle & オラクル & テスト話者で最良の値．選び方の損を測るためだけに使う \\
    student / teacher & 生徒／先生 & 学ぶ線形分類器／弱いラベルを出す分類器 \\
    weak labels & 弱いラベル & 先生がラベル無しの発話に付ける感情の確率（感情ごとの合計をそろえる） \\
    fused start & 融合の出発点 & 融合した特徴でソースだけから作った出発点（基本版の先生の候補の 1 つ） \\
    \bottomrule
  \end{tabularx}
\end{frame}

\begin{frame}{付録 1：論文の英語の語と，この資料の語（4／4：比較）}
  \rl{語の対応}
  \vspace{-0.2em}
  \centering\scriptsize
  \renewcommand{\arraystretch}{1.12}
  \begin{tabularx}{\textwidth}{@{}l l X@{}}
    \toprule
    論文の語 & この資料の語 & 意味 \\
    \midrule
    two-stage reference & 融合の出発点が教えた単独の生徒の予測の積 & 基本版の融合の生徒の先生の候補 \\
    belief & 見込み & ソースごとの分類器を混同の仕方で重みづけてまとめた先生（同じモダリティ） \\
    $J$ & 先生の優位 $J$ & プールでの先生の UAR − 出発点の UAR \\
    ours / ours-ext & 基本版／拡張版 & 拡張版は $K$ 本で学び直した融合の生徒を先生にし，重みづけも選ぶ \\
    tuned / fixed & 調整版／固定版 & 設定を仮のターゲットで選ぶ／元の論文の値を使う \\
    strongest tuned competitor & 最も強い比較手法 & 調整版の比較手法のうち，ターゲットの平均が最も高いもの \\
    rebalancing / column normalization & 釣り合い直し／列の割り直し & 感情ごとの合計をそろえる操作 \\
    \bottomrule
  \end{tabularx}
\end{frame}

\begin{frame}{付録 2：記号の一覧（1／2：第 I 部）}
  \rl{記号}\label{a:symbols}
  \vspace{-0.2em}
  \centering\scriptsize
  \renewcommand{\arraystretch}{1.15}
  \begin{tabularx}{\textwidth}{@{}l X X@{}}
    \toprule
    記号 & 何か & 大きさの目安 \\
    \midrule
    $K$ & 感情あたりのターゲットのラベルの本数 & 0〜50 と全件 \\
    $z$ & 2 つの感情のスコアの差（線形分類器の出力の差） & -- \\
    $d$，$\Delta$ & 識別力（今の向き），ターゲットの全ラベルで求めた最良の向きの識別力 & 2 前後で正解率 84\% \\
    $o$ & 判断基準のずれ（山の幅で割る） & 1 で正解率が約 10 \pt 下がる（$d=2$） \\
    $a(d,o)$ & 2 つの感情の正解率（2 つの再現率の平均） & 50〜100\% \\
    $q$ & 向きの誤差の減り方を決める数 & A \NUM{p1-qdir-a}，A(HuBERT) \NUM{p1-qdir-ah}，AV \NUM{p1-qdir-av}，V \NUM{p1-qdir-v} \\
    $e$ & ソースの向きの効率（$d = e\Delta$） & 0〜1 \\
    $n$ & ソースの価値（感情あたりのラベルの本数） & 音声の 1 ソースで \NUM{p1-e32-n1src-audio-lo}〜\NUM{p1-e32-n1src-audio-hi} \\
    $\lambda$，$\rho$ & 合併が判断基準のずれを直す割合，それを決める量（$\lambda = \rho/(1+\rho)$） & $\lambda$ は 0〜1 \\
    \bottomrule
  \end{tabularx}
\end{frame}

\begin{frame}{付録 2：記号の一覧（2／2：第 II 部）}
  \rl{記号}
  \vspace{-0.2em}
  \centering\scriptsize
  \renewcommand{\arraystretch}{1.15}
  \begin{tabularx}{\textwidth}{@{}l X X@{}}
    \toprule
    記号 & 何か & 大きさの目安 \\
    \midrule
    $W_0, b_0$ & 出発点の重みと切片 & -- \\
    $W^{*}$，$E_0$ & プールの全ラベルで学んだ生徒，出発点の向きの誤差 $W^{*} - W_0$ & -- \\
    $C'$，$C_b$，$\rho_M$，$\gamma$，$\rho_s$ & 5 つの値（重みと切片の縛りの強さ，計量の混ぜる比，ソースの比重の指数，ソースの向きの縮め方） & 3 つの候補から仮のターゲットで選ぶ \\
    $Q$，$m$ & 先生の予測（発話 × 4 感情），弱いラベル全体の比重 & $m$ は $K$ ごとに選ぶ \\
    $J$ & 先生の優位（プールでの UAR の差） & 見込み \NUM{p2-e52-q-minus-p-a}，融合の出発点 \NUM{p2-e60pre-juar-av-a-fused}（A，\pt） \\
    \bottomrule
  \end{tabularx}
\end{frame}

\begin{frame}{付録 3：研究の記録の実験の番号と，台本の中の記号}
  \rl{実験の番号と記号}\label{a:codes}
  \vspace{-0.2em}
  \begin{columns}[T,onlytextwidth]
    \begin{column}{0.56\textwidth}
      \centering\tiny
      \renewcommand{\arraystretch}{1.08}
      \begin{tabularx}{\linewidth}{@{}l X@{}}
        \toprule
        番号 & この資料での中身 \\
        \midrule
        E12 & 中心化と，ソースのみの切片だけを学び直す（D1） \\
        E13 & 損を判断基準による損失と向きによる損失に分ける \\
        E15，E16 & $\lambda$ の式と，Newton 法の 1 歩の全体 \\
        E28 & $q$ のあてはめ（E28-9）など手法の数理の確かめ \\
        E29，E30 & 判断基準の割合（線形，2 層のパーセプトロン） \\
        E31 & 全ソースからの仮説転移（P4）と従来法の比較 \\
        E32 & ソースの価値 $n$ \\
        E39 & 5 つの値を選んだ出発点 \\
        E49，E50，E53 & 5 つの値の生徒でのラベル無しの選び方 \\
        E52，E59 & 見込みと見込みの蒸留 \\
        E58 & 出発点の向きの誤差 $E_0$ を捉える割合 \\
        E60 & モダリティをまたぐ蒸留（基本版） \\
        E61，E62，E63，E64 & 比較手法，2 つ目のエンコーダ，微調整，検定 \\
        \bottomrule
      \end{tabularx}
    \end{column}
    \begin{column}{0.42\textwidth}
      \centering\tiny
      \renewcommand{\arraystretch}{1.08}
      \begin{tabularx}{\linewidth}{@{}l X@{}}
        \toprule
        記号 & この資料の語 \\
        \midrule
        e39 & 5 つの値を選んだ出発点 \\
        P4 & 全ソースからの仮説転移 \\
        D1（E12） & ソースのみの切片だけを $K$ 本で学び直す \\
        D1（E39） & ソースを等しい比重で足す出発点 \\
        B1，B4 & ターゲットのみ，合併（E31） \\
        pF & 融合の出発点 \\
        pFK & $K$ 本で学び直した融合の生徒 \\
        ref & 融合の出発点が教えた単独の生徒の予測の積 \\
        ours，oursx & 基本版，拡張版 \\
        meta & 仮のターゲットで選んだ版 \\
        W7，Hu7 & A，A(HuBERT) \\
        \bottomrule
      \end{tabularx}
    \end{column}
  \end{columns}
  \vspace{0.3em}
  {\raggedright\tiny D1 は 2 つの実験で別のものを指す．本文と補足では記号を使わず，中身で書いた\par}
\end{frame}

\IfFileExists{figures/appendix_sources.tex}{% 生成物（tools/key_sources.py）．資料の \NUM の key 761 個を 111 の頭にまとめた．手で直さない
\begin{frame}{付録 4：数値の出所（1／12）}
  \rl{数値の出所}\label{a:sources}
  \vspace{-0.3em}\centering\tiny\renewcommand{\arraystretch}{1.0}
  \begin{tabularx}{\textwidth}{@{}l r l X@{}}
    \toprule key の頭 & 数 & 使った part & 台帳の出所（結果ファイル） \\ \midrule
    \texttt{cmp-best-vs} & 22 & m\_compare & \texttt{results/\allowbreak{}baselines\_\allowbreak{}fuse-audio-text/\allowbreak{}<ターゲット>.csv；results/\allowbreak{}baselines\_\allowbreak{}fuse-audio-video/\allowbreak{}<ターゲット>.csv；ほか} \\
    \texttt{cmp-coreg-saerens} & 2 & supp\_compare & \texttt{results/\allowbreak{}baselines\_\allowbreak{}<世界>/\allowbreak{}<ターゲット>.csv} \\
    \texttt{cmp-entmin-fixed} & 34 & m\_compare, supp\_compare & \texttt{results/\allowbreak{}baselines\_\allowbreak{}fuse-audio-text/\allowbreak{}<ターゲット>.csv；results/\allowbreak{}baselines\_\allowbreak{}fuse-audio-text/\allowbreak{}IEMOCAP-impro.csv；ほか} \\
    \texttt{cmp-fixed-collapse} & 2 & supp\_compare & \texttt{results/\allowbreak{}baselines\_\allowbreak{}fuse-audio-video-text/\allowbreak{}<ターゲット>.csv} \\
    \texttt{cmp-frozen-uar} & 4 & supp\_compare & \texttt{results/\allowbreak{}finetune\_\allowbreak{}check/\allowbreak{}teacher\_\allowbreak{}<ターゲット>.csv} \\
    \texttt{cmp-ft-anchor} & 3 & supp\_compare & \texttt{results/\allowbreak{}finetune\_\allowbreak{}check/\allowbreak{}<ターゲット>.csv} \\
    \texttt{cmp-ft-konly} & 3 & m\_compare, supp\_compare & \texttt{results/\allowbreak{}finetune\_\allowbreak{}check/\allowbreak{}<ターゲット>.csv；results/\allowbreak{}finetune\_\allowbreak{}check/\allowbreak{}teacher\_\allowbreak{}<ターゲット>.csv} \\
    \texttt{cmp-ft-ours} & 6 & m\_compare, supp\_compare & \texttt{results/\allowbreak{}finetune\_\allowbreak{}check/\allowbreak{}<ターゲット>.csv；results/\allowbreak{}finetune\_\allowbreak{}check/\allowbreak{}teacher\_\allowbreak{}<ターゲット>.csv} \\
    \texttt{cmp-ft-pairs} & 1 & supp\_compare & \texttt{results/\allowbreak{}finetune\_\allowbreak{}check/\allowbreak{}<ターゲット>.csv} \\
    \texttt{cmp-ft-reps} & 1 & supp\_compare & \texttt{results/\allowbreak{}finetune\_\allowbreak{}check/\allowbreak{}<ターゲット>.csv；results/\allowbreak{}finetune\_\allowbreak{}check/\allowbreak{}<ターゲット>\_\allowbreak{}c.csv；ほか} \\
    \bottomrule
  \end{tabularx}\par
  {\raggedright\tiny key の実体（値・書式・出所・確かめた日）は paper/results/numbers.json．数は資料で使った key の数\par}
\end{frame}
\begin{frame}{付録 4：数値の出所（2／12）}
  \rl{数値の出所}
  \vspace{-0.3em}\centering\tiny\renewcommand{\arraystretch}{1.0}
  \begin{tabularx}{\textwidth}{@{}l r l X@{}}
    \toprule key の頭 & 数 & 使った part & 台帳の出所（結果ファイル） \\ \midrule
    \texttt{cmp-ft-start} & 1 & supp\_compare & \texttt{results/\allowbreak{}finetune\_\allowbreak{}check/\allowbreak{}<ターゲット>.csv；results/\allowbreak{}finetune\_\allowbreak{}check/\allowbreak{}teacher\_\allowbreak{}<ターゲット>.csv} \\
    \texttt{cmp-ft-uar} & 5 & supp\_compare & \texttt{results/\allowbreak{}finetune\_\allowbreak{}check/\allowbreak{}<ターゲット>.csv} \\
    \texttt{cmp-ft-weak} & 4 & m\_compare, supp\_compare & \texttt{results/\allowbreak{}finetune\_\allowbreak{}check/\allowbreak{}<ターゲット>.csv} \\
    \texttt{cmp-ftc-anchor} & 1 & supp\_compare & \texttt{results/\allowbreak{}finetune\_\allowbreak{}check/\allowbreak{}<ターゲット>\_\allowbreak{}c.csv} \\
    \texttt{cmp-ftc-gap} & 2 & supp\_compare & \texttt{results/\allowbreak{}finetune\_\allowbreak{}check/\allowbreak{}<ターゲット>\_\allowbreak{}c.csv；results/\allowbreak{}finetune\_\allowbreak{}check/\allowbreak{}<ターゲット>.csv；ほか} \\
    \texttt{cmp-ftc-konly} & 3 & supp\_compare & \texttt{results/\allowbreak{}finetune\_\allowbreak{}check/\allowbreak{}<ターゲット>\_\allowbreak{}c.csv；results/\allowbreak{}finetune\_\allowbreak{}check/\allowbreak{}teacher\_\allowbreak{}<ターゲット>.csv} \\
    \texttt{cmp-ftc-lift} & 4 & supp\_compare & \texttt{results/\allowbreak{}finetune\_\allowbreak{}check/\allowbreak{}<ターゲット>\_\allowbreak{}c.csv；results/\allowbreak{}finetune\_\allowbreak{}check/\allowbreak{}<ターゲット>.csv} \\
    \texttt{cmp-ftc-ours} & 6 & supp\_compare & \texttt{results/\allowbreak{}finetune\_\allowbreak{}check/\allowbreak{}<ターゲット>\_\allowbreak{}c.csv；results/\allowbreak{}finetune\_\allowbreak{}check/\allowbreak{}teacher\_\allowbreak{}<ターゲット>.csv} \\
    \texttt{cmp-ftc-start} & 1 & supp\_compare & \texttt{results/\allowbreak{}finetune\_\allowbreak{}check/\allowbreak{}<ターゲット>\_\allowbreak{}c.csv；results/\allowbreak{}finetune\_\allowbreak{}check/\allowbreak{}teacher\_\allowbreak{}<ターゲット>.csv} \\
    \texttt{cmp-ftc-uar} & 5 & supp\_compare & \texttt{results/\allowbreak{}finetune\_\allowbreak{}check/\allowbreak{}<ターゲット>\_\allowbreak{}c.csv} \\
    \bottomrule
  \end{tabularx}\par
  {\raggedright\tiny key の実体（値・書式・出所・確かめた日）は paper/results/numbers.json．数は資料で使った key の数\par}
\end{frame}
\begin{frame}{付録 4：数値の出所（3／12）}
  \rl{数値の出所}
  \vspace{-0.3em}\centering\tiny\renewcommand{\arraystretch}{1.0}
  \begin{tabularx}{\textwidth}{@{}l r l X@{}}
    \toprule key の頭 & 数 & 使った part & 台帳の出所（結果ファイル） \\ \midrule
    \texttt{cmp-ftc-weak} & 4 & supp\_compare & \texttt{results/\allowbreak{}finetune\_\allowbreak{}check/\allowbreak{}<ターゲット>\_\allowbreak{}c.csv} \\
    \texttt{cmp-gd-vs} & 2 & m\_compare, supp\_compare & \texttt{results/\allowbreak{}baselines\_\allowbreak{}fuse-audio-video-text/\allowbreak{}<ターゲット>.csv} \\
    \texttt{cmp-ours-vs} & 27 & m\_compare, m\_intro, supp\_compare & \texttt{results/\allowbreak{}baselines\_\allowbreak{}fuse-audio-text/\allowbreak{}<ターゲット>.csv；results/\allowbreak{}multimodal\_\allowbreak{}distill\_\allowbreak{}fuse-audio-text/\allowbreak{}<ターゲット>.csv；ほか} \\
    \texttt{cmp-oursx-oracle} & 13 & m\_compare, m\_end, supp\_compare & \texttt{paper/\allowbreak{}scripts/\allowbreak{}ledger/\allowbreak{}comparison.py；results/\allowbreak{}baselines\_\allowbreak{}fuse-audio-text/\allowbreak{}<ターゲット>.csv；ほか} \\
    \texttt{cmp-oursx-vs} & 42 & m\_compare, supp\_compare & \texttt{results/\allowbreak{}baselines\_\allowbreak{}fuse-audio-video-text/\allowbreak{}<ターゲット>.csv；paper/\allowbreak{}scripts/\allowbreak{}ledger/\allowbreak{}comparison.py；ほか} \\
    \texttt{cmp-pseudo-reps} & 1 & supp\_compare, supp\_setting & \texttt{results/\allowbreak{}baselines\_\allowbreak{}<世界>/\allowbreak{}<ターゲット>\_\allowbreak{}pseudo.csv} \\
    \texttt{cmp-reps-n} & 1 & m\_compare, supp\_compare & \texttt{results/\allowbreak{}baselines\_\allowbreak{}<世界>/\allowbreak{}<ターゲット>.csv；results/\allowbreak{}multimodal\_\allowbreak{}distill\_\allowbreak{}<世界>/\allowbreak{}<ターゲット>.csv} \\
    \texttt{cmp-targets-at} & 1 & supp\_compare & \texttt{results/\allowbreak{}baselines\_\allowbreak{}fuse-audio-text/\allowbreak{}} \\
    \texttt{cmp-targets-av} & 1 & supp\_compare & \texttt{results/\allowbreak{}baselines\_\allowbreak{}fuse-audio-video/\allowbreak{}} \\
    \texttt{cmp-targets-avt} & 1 & supp\_compare & \texttt{results/\allowbreak{}baselines\_\allowbreak{}fuse-audio-video-text/\allowbreak{}} \\
    \bottomrule
  \end{tabularx}\par
  {\raggedright\tiny key の実体（値・書式・出所・確かめた日）は paper/results/numbers.json．数は資料で使った key の数\par}
\end{frame}
\begin{frame}{付録 4：数値の出所（4／12）}
  \rl{数値の出所}
  \vspace{-0.3em}\centering\tiny\renewcommand{\arraystretch}{1.0}
  \begin{tabularx}{\textwidth}{@{}l r l X@{}}
    \toprule key の頭 & 数 & 使った part & 台帳の出所（結果ファイル） \\ \midrule
    \texttt{cmp-targets-vt} & 1 & supp\_compare & \texttt{results/\allowbreak{}baselines\_\allowbreak{}fuse-video-text/\allowbreak{}} \\
    \texttt{cmp-targets-w2} & 1 & supp\_compare & \texttt{results/\allowbreak{}baselines\_\allowbreak{}fuse2-audio-video-text/\allowbreak{}} \\
    \texttt{data-afew} & 4 & m\_setting & \texttt{measure/\allowbreak{}data.py load\_\allowbreak{}group('AFEW', 'audio\_\allowbreak{}wavlm')} \\
    \texttt{data-cremad} & 4 & m\_setting & \texttt{measure/\allowbreak{}data.py load\_\allowbreak{}group('CREMA-D', 'audio\_\allowbreak{}wavlm')} \\
    \texttt{data-impro} & 4 & m\_setting & \texttt{measure/\allowbreak{}data.py load\_\allowbreak{}group('IEMOCAP-impro', 'audio\_\allowbreak{}wavlm')} \\
    \texttt{data-meld} & 4 & m\_setting & \texttt{measure/\allowbreak{}data.py load\_\allowbreak{}group('MELD', 'audio\_\allowbreak{}wavlm')} \\
    \texttt{data-mosei} & 7 & m\_setting, supp\_setting & \texttt{measure/\allowbreak{}data.py load\_\allowbreak{}group('MOSEI', 'audio\_\allowbreak{}wavlm')；data/\allowbreak{}metadata/\allowbreak{}mosei\_\allowbreak{}channel\_\allowbreak{}map.csv, data/\allowbreak{}metadata/\allowbreak{}mosei\_\allowbreak{}utterance\_\allowbreak{}manifest.csv；ほか} \\
    \texttt{data-n} & 47 & m\_setting, supp\_setting & \texttt{measure/\allowbreak{}data.py load\_\allowbreak{}any('AFEW', 'fuse-audio-video')；measure/\allowbreak{}data.py load\_\allowbreak{}group('AFEW', 'video\_\allowbreak{}clip')；ほか} \\
    \texttt{data-ravdess} & 4 & m\_setting & \texttt{measure/\allowbreak{}data.py load\_\allowbreak{}group('RAVDESS', 'audio\_\allowbreak{}wavlm')} \\
    \texttt{data-script} & 4 & m\_setting & \texttt{measure/\allowbreak{}data.py load\_\allowbreak{}group('IEMOCAP-script', 'audio\_\allowbreak{}wavlm')} \\
    \bottomrule
  \end{tabularx}\par
  {\raggedright\tiny key の実体（値・書式・出所・確かめた日）は paper/results/numbers.json．数は資料で使った key の数\par}
\end{frame}
\begin{frame}{付録 4：数値の出所（5／12）}
  \rl{数値の出所}
  \vspace{-0.3em}\centering\tiny\renewcommand{\arraystretch}{1.0}
  \begin{tabularx}{\textwidth}{@{}l r l X@{}}
    \toprule key の頭 & 数 & 使った part & 台帳の出所（結果ファイル） \\ \midrule
    \texttt{data-spk} & 7 & m\_setting & \texttt{measure/\allowbreak{}data.py load\_\allowbreak{}group('AFEW', 'audio\_\allowbreak{}wavlm')；measure/\allowbreak{}data.py load\_\allowbreak{}group('CREMA-D', 'audio\_\allowbreak{}wavlm')；ほか} \\
    \texttt{data-test} & 3 & supp\_setting & \texttt{results/\allowbreak{}fixc\_\allowbreak{}map/\allowbreak{}audio\_\allowbreak{}wavlm/\allowbreak{}*.csv, measure/\allowbreak{}data.py load\_\allowbreak{}group(群, 'audio\_\allowbreak{}wavlm')} \\
    \texttt{map-gain} & 3 & m\_part1 & \texttt{results/\allowbreak{}fixc\_\allowbreak{}map/\allowbreak{}audio\_\allowbreak{}wavlm/\allowbreak{}summary/\allowbreak{}pairs.csv} \\
    \texttt{map-pairs} & 16 & m\_part1, supp\_part1 & \texttt{results/\allowbreak{}fixc\_\allowbreak{}map/\allowbreak{}audio\_\allowbreak{}wavlm/\allowbreak{}*.csv；results/\allowbreak{}fixc\_\allowbreak{}map/\allowbreak{}audio\_\allowbreak{}hubert/\allowbreak{}*.csv；ほか} \\
    \texttt{map-reps} & 1 & m\_setting, supp\_part1 & \texttt{results/\allowbreak{}fixc\_\allowbreak{}map/\allowbreak{}audio\_\allowbreak{}wavlm/\allowbreak{}*.csv} \\
    \texttt{map-share} & 2 & m\_part1, supp\_part1 & \texttt{results/\allowbreak{}fixc\_\allowbreak{}analysis/\allowbreak{}explore\_\allowbreak{}relation/\allowbreak{}snapshots/\allowbreak{}20261009-0105\_\allowbreak{}fixc/\allowbreak{}decomp\_\allowbreak{}share.csv；results/\allowbreak{}rerun\_\allowbreak{}fixc/\allowbreak{}check\_\allowbreak{}h7/\allowbreak{}share.csv} \\
    \texttt{p1-center-d1} & 3 & m\_part1, supp\_part1 & \texttt{results/\allowbreak{}mech\_\allowbreak{}audio\_\allowbreak{}wavlm/\allowbreak{}*.csv；results/\allowbreak{}mech\_\allowbreak{}audio\_\allowbreak{}hubert/\allowbreak{}*.csv；ほか} \\
    \texttt{p1-center-merged} & 3 & m\_part1, supp\_part1 & \texttt{results/\allowbreak{}mech\_\allowbreak{}audio\_\allowbreak{}wavlm/\allowbreak{}*.csv；results/\allowbreak{}mech\_\allowbreak{}audio\_\allowbreak{}hubert/\allowbreak{}*.csv；ほか} \\
    \texttt{p1-center-srconly} & 9 & m\_part1, supp\_part1 & \texttt{results/\allowbreak{}mech\_\allowbreak{}audio\_\allowbreak{}wavlm/\allowbreak{}*.csv；results/\allowbreak{}mech\_\allowbreak{}audio\_\allowbreak{}hubert/\allowbreak{}*.csv；ほか} \\
    \texttt{p1-d1-gain} & 3 & m\_part1, supp\_part1 & \texttt{results/\allowbreak{}mech\_\allowbreak{}audio\_\allowbreak{}wavlm/\allowbreak{}*.csv；results/\allowbreak{}mech\_\allowbreak{}audio\_\allowbreak{}hubert/\allowbreak{}*.csv；ほか} \\
    \bottomrule
  \end{tabularx}\par
  {\raggedright\tiny key の実体（値・書式・出所・確かめた日）は paper/results/numbers.json．数は資料で使った key の数\par}
\end{frame}
\begin{frame}{付録 4：数値の出所（6／12）}
  \rl{数値の出所}
  \vspace{-0.3em}\centering\tiny\renewcommand{\arraystretch}{1.0}
  \begin{tabularx}{\textwidth}{@{}l r l X@{}}
    \toprule key の頭 & 数 & 使った part & 台帳の出所（結果ファイル） \\ \midrule
    \texttt{p1-dirloss-recov} & 10 & m\_part1, supp\_part1 & \texttt{results/\allowbreak{}explore\_\allowbreak{}relation/\allowbreak{}snapshots/\allowbreak{}20261011-0139\_\allowbreak{}e28t9/\allowbreak{}e28t9\_\allowbreak{}pairs.csv；results/\allowbreak{}explore\_\allowbreak{}relation/\allowbreak{}snapshots/\allowbreak{}20261011-0139\_\allowbreak{}e28t9/\allowbreak{}e28t9\_\allowbreak{}pairs.csv・e28t9\_\allowbreak{}targets.csv} \\
    \texttt{p1-dirloss-remain} & 4 & supp\_part1 & \texttt{results/\allowbreak{}explore\_\allowbreak{}relation/\allowbreak{}snapshots/\allowbreak{}20261011-0139\_\allowbreak{}e28t9/\allowbreak{}e28t9\_\allowbreak{}pairs.csv} \\
    \texttt{p1-e12-reps} & 1 & supp\_part1 & \texttt{results/\allowbreak{}mech\_\allowbreak{}{audio\_\allowbreak{}wavlm,audio\_\allowbreak{}hubert,video}/\allowbreak{}*.csv} \\
    \texttt{p1-e13-fit} & 18 & m\_setting, supp\_part1 & \texttt{results/\allowbreak{}rerun\_\allowbreak{}fixc/\allowbreak{}explore\_\allowbreak{}relation/\allowbreak{}snapshots/\allowbreak{}20261018-0213\_\allowbreak{}fixc/\allowbreak{}e13\_\allowbreak{}pairs.csv；results/\allowbreak{}rerun\_\allowbreak{}fixc/\allowbreak{}explore\_\allowbreak{}relation/\allowbreak{}snapshots/\allowbreak{}20261018-0213\_\allowbreak{}fixc/\allowbreak{}check\_\allowbreak{}mh/\allowbreak{}e13.csv} \\
    \texttt{p1-e13-loss} & 30 & m\_part1, supp\_compare, supp\_part1 & \texttt{results/\allowbreak{}rerun\_\allowbreak{}fixc/\allowbreak{}explore\_\allowbreak{}relation/\allowbreak{}snapshots/\allowbreak{}20261018-0213\_\allowbreak{}fixc/\allowbreak{}e13\_\allowbreak{}pairs.csv；results/\allowbreak{}rerun\_\allowbreak{}fixc/\allowbreak{}explore\_\allowbreak{}relation/\allowbreak{}snapshots/\allowbreak{}20261018-0213\_\allowbreak{}fixc/\allowbreak{}check\_\allowbreak{}mh/\allowbreak{}e13.csv} \\
    \texttt{p1-e15-lam} & 6 & m\_part1, supp\_part1 & \texttt{results/\allowbreak{}fixc\_\allowbreak{}analysis/\allowbreak{}explore\_\allowbreak{}relation/\allowbreak{}snapshots/\allowbreak{}20261009-0105\_\allowbreak{}fixc/\allowbreak{}stageB/\allowbreak{}e15b\_\allowbreak{}pairs.csv} \\
    \texttt{p1-e16-dir} & 4 & supp\_part1 & \texttt{results/\allowbreak{}fixc\_\allowbreak{}analysis/\allowbreak{}explore\_\allowbreak{}relation/\allowbreak{}snapshots/\allowbreak{}20261009-0105\_\allowbreak{}fixc/\allowbreak{}stageB/\allowbreak{}e16\_\allowbreak{}pairs.csv} \\
    \texttt{p1-e16-lam} & 4 & supp\_part1 & \texttt{results/\allowbreak{}fixc\_\allowbreak{}analysis/\allowbreak{}explore\_\allowbreak{}relation/\allowbreak{}snapshots/\allowbreak{}20261009-0105\_\allowbreak{}fixc/\allowbreak{}stageB/\allowbreak{}e16\_\allowbreak{}pairs.csv} \\
    \texttt{p1-e28t9-reps} & 1 & supp\_part1 & \texttt{results/\allowbreak{}explore\_\allowbreak{}relation/\allowbreak{}snapshots/\allowbreak{}20261011-0139\_\allowbreak{}e28t9/\allowbreak{}provenance\_\allowbreak{}audio\_\allowbreak{}wavlm.json} \\
    \texttt{p1-e29-critcorr} & 3 & supp\_part1 & \texttt{results/\allowbreak{}explore\_\allowbreak{}relation/\allowbreak{}snapshots/\allowbreak{}20261011-1524\_\allowbreak{}e29/\allowbreak{}e29\_\allowbreak{}pairs.csv} \\
    \bottomrule
  \end{tabularx}\par
  {\raggedright\tiny key の実体（値・書式・出所・確かめた日）は paper/results/numbers.json．数は資料で使った key の数\par}
\end{frame}
\begin{frame}{付録 4：数値の出所（7／12）}
  \rl{数値の出所}
  \vspace{-0.3em}\centering\tiny\renewcommand{\arraystretch}{1.0}
  \begin{tabularx}{\textwidth}{@{}l r l X@{}}
    \toprule key の頭 & 数 & 使った part & 台帳の出所（結果ファイル） \\ \midrule
    \texttt{p1-e29-critshare} & 28 & m\_end, m\_intro, m\_part1, supp\_compare, supp\_part1 & \texttt{results/\allowbreak{}explore\_\allowbreak{}relation/\allowbreak{}snapshots/\allowbreak{}20261011-1524\_\allowbreak{}e29/\allowbreak{}e29\_\allowbreak{}pairs.csv} \\
    \texttt{p1-e29-lossshare} & 10 & m\_part1, supp\_part1 & \texttt{results/\allowbreak{}explore\_\allowbreak{}relation/\allowbreak{}snapshots/\allowbreak{}20261011-1524\_\allowbreak{}e29/\allowbreak{}e29\_\allowbreak{}pairs.csv} \\
    \texttt{p1-e29-shares} & 1 & supp\_part1 & \texttt{results/\allowbreak{}explore\_\allowbreak{}relation/\allowbreak{}snapshots/\allowbreak{}20261011-1524\_\allowbreak{}e29/\allowbreak{}e29\_\allowbreak{}pairs.csv} \\
    \texttt{p1-e30-critshare} & 8 & supp\_part1 & \texttt{results/\allowbreak{}explore\_\allowbreak{}relation/\allowbreak{}snapshots/\allowbreak{}20261013-1346\_\allowbreak{}e30/\allowbreak{}e30\_\allowbreak{}pairs.csv} \\
    \texttt{p1-e30-reps} & 1 & supp\_part1 & \texttt{results/\allowbreak{}explore\_\allowbreak{}relation/\allowbreak{}snapshots/\allowbreak{}20261013-1346\_\allowbreak{}e30/\allowbreak{}provenance\_\allowbreak{}audio\_\allowbreak{}wavlm.json} \\
    \texttt{p1-e32-allsrc} & 8 & supp\_part1 & \texttt{results/\allowbreak{}explore\_\allowbreak{}relation/\allowbreak{}snapshots/\allowbreak{}20261016-0954\_\allowbreak{}e32/\allowbreak{}e32\_\allowbreak{}allsrc.csv} \\
    \texttt{p1-e32-n1src} & 14 & appendix, m\_end, m\_intro, m\_part1, supp\_part1 & \texttt{results/\allowbreak{}explore\_\allowbreak{}relation/\allowbreak{}snapshots/\allowbreak{}20261016-0954\_\allowbreak{}e32/\allowbreak{}e32\_\allowbreak{}pairs.csv・e32\_\allowbreak{}allsrc.csv；results/\allowbreak{}explore\_\allowbreak{}relation/\allowbreak{}snapshots/\allowbreak{}20261016-0954\_\allowbreak{}e32/\allowbreak{}e32\_\allowbreak{}pairs.csv} \\
    \texttt{p1-e32-nall} & 10 & m\_intro, m\_part1, supp\_part1 & \texttt{results/\allowbreak{}explore\_\allowbreak{}relation/\allowbreak{}snapshots/\allowbreak{}20261016-0954\_\allowbreak{}e32/\allowbreak{}e32\_\allowbreak{}pairs.csv・e32\_\allowbreak{}allsrc.csv；results/\allowbreak{}explore\_\allowbreak{}relation/\allowbreak{}snapshots/\allowbreak{}20261016-0954\_\allowbreak{}e32/\allowbreak{}e32\_\allowbreak{}allsrc.csv} \\
    \texttt{p1-e32-ratio} & 8 & supp\_part1 & \texttt{results/\allowbreak{}explore\_\allowbreak{}relation/\allowbreak{}snapshots/\allowbreak{}20261016-0954\_\allowbreak{}e32/\allowbreak{}e32\_\allowbreak{}pairs.csv} \\
    \texttt{p1-e32-rho} & 10 & m\_part1, supp\_part1 & \texttt{results/\allowbreak{}explore\_\allowbreak{}relation/\allowbreak{}snapshots/\allowbreak{}20261016-0954\_\allowbreak{}e32/\allowbreak{}e32\_\allowbreak{}pairs.csv} \\
    \bottomrule
  \end{tabularx}\par
  {\raggedright\tiny key の実体（値・書式・出所・確かめた日）は paper/results/numbers.json．数は資料で使った key の数\par}
\end{frame}
\begin{frame}{付録 4：数値の出所（8／12）}
  \rl{数値の出所}
  \vspace{-0.3em}\centering\tiny\renewcommand{\arraystretch}{1.0}
  \begin{tabularx}{\textwidth}{@{}l r l X@{}}
    \toprule key の頭 & 数 & 使った part & 台帳の出所（結果ファイル） \\ \midrule
    \texttt{p1-merged-gain} & 3 & m\_part1, supp\_part1 & \texttt{results/\allowbreak{}mech\_\allowbreak{}audio\_\allowbreak{}wavlm/\allowbreak{}*.csv；results/\allowbreak{}mech\_\allowbreak{}audio\_\allowbreak{}hubert/\allowbreak{}*.csv；ほか} \\
    \texttt{p1-qdir-a} & 1 & appendix, m\_part1, supp\_part1 & \texttt{results/\allowbreak{}explore\_\allowbreak{}relation/\allowbreak{}snapshots/\allowbreak{}20261011-0139\_\allowbreak{}e28t9/\allowbreak{}e28t9\_\allowbreak{}targets.csv} \\
    \texttt{p1-qdir-ah} & 1 & appendix, m\_part1, supp\_part1 & \texttt{results/\allowbreak{}explore\_\allowbreak{}relation/\allowbreak{}snapshots/\allowbreak{}20261011-0139\_\allowbreak{}e28t9/\allowbreak{}e28t9\_\allowbreak{}targets.csv} \\
    \texttt{p1-qdir-audio} & 2 & m\_intro & \texttt{results/\allowbreak{}explore\_\allowbreak{}relation/\allowbreak{}snapshots/\allowbreak{}20261011-0139\_\allowbreak{}e28t9/\allowbreak{}e28t9\_\allowbreak{}pairs.csv・e28t9\_\allowbreak{}targets.csv} \\
    \texttt{p1-qdir-av} & 1 & appendix, m\_part1, supp\_part1 & \texttt{results/\allowbreak{}explore\_\allowbreak{}relation/\allowbreak{}snapshots/\allowbreak{}20261011-0139\_\allowbreak{}e28t9/\allowbreak{}e28t9\_\allowbreak{}targets.csv} \\
    \texttt{p1-qdir-v} & 1 & appendix, m\_part1, supp\_part1 & \texttt{results/\allowbreak{}explore\_\allowbreak{}relation/\allowbreak{}snapshots/\allowbreak{}20261011-0139\_\allowbreak{}e28t9/\allowbreak{}e28t9\_\allowbreak{}targets.csv} \\
    \texttt{p1-stageb-reps} & 1 & supp\_part1 & \texttt{results/\allowbreak{}fixc\_\allowbreak{}analysis/\allowbreak{}explore\_\allowbreak{}relation/\allowbreak{}snapshots/\allowbreak{}20261009-0105\_\allowbreak{}fixc/\allowbreak{}stageB/\allowbreak{}provenance\_\allowbreak{}merged\_\allowbreak{}scores\_\allowbreak{}audio\_\allowbreak{}wavlm.json} \\
    \texttt{p2-cond2-e39} & 6 & supp\_part2 & \texttt{results/\allowbreak{}explore\_\allowbreak{}relation/\allowbreak{}snapshots/\allowbreak{}20261021-0440\_\allowbreak{}e49/\allowbreak{}e49\_\allowbreak{}targets.csv；results/\allowbreak{}explore\_\allowbreak{}relation/\allowbreak{}snapshots/\allowbreak{}20261021-1220\_\allowbreak{}e50/\allowbreak{}e50\_\allowbreak{}targets.csv；ほか} \\
    \texttt{p2-cond2-student} & 7 & supp\_part2 & \texttt{results/\allowbreak{}explore\_\allowbreak{}relation/\allowbreak{}snapshots/\allowbreak{}20261021-0440\_\allowbreak{}e49/\allowbreak{}e49\_\allowbreak{}targets.csv；results/\allowbreak{}explore\_\allowbreak{}relation/\allowbreak{}snapshots/\allowbreak{}20261021-1450\_\allowbreak{}e53/\allowbreak{}e53\_\allowbreak{}targets.csv} \\
    \texttt{p2-e31-p4} & 12 & m\_part2, supp\_part2 & \texttt{results/\allowbreak{}explore\_\allowbreak{}relation/\allowbreak{}snapshots/\allowbreak{}20261017-1654\_\allowbreak{}e31final/\allowbreak{}e31\_\allowbreak{}finals30.csv} \\
    \bottomrule
  \end{tabularx}\par
  {\raggedright\tiny key の実体（値・書式・出所・確かめた日）は paper/results/numbers.json．数は資料で使った key の数\par}
\end{frame}
\begin{frame}{付録 4：数値の出所（9／12）}
  \rl{数値の出所}
  \vspace{-0.3em}\centering\tiny\renewcommand{\arraystretch}{1.0}
  \begin{tabularx}{\textwidth}{@{}l r l X@{}}
    \toprule key の頭 & 数 & 使った part & 台帳の出所（結果ファイル） \\ \midrule
    \texttt{p2-e39-vs} & 10 & m\_part2, supp\_part2 & \texttt{results/\allowbreak{}explore\_\allowbreak{}relation/\allowbreak{}snapshots/\allowbreak{}20261018-0213\_\allowbreak{}ext30/\allowbreak{}ext30\_\allowbreak{}targets.csv} \\
    \texttt{p2-e52-q} & 2 & appendix, m\_part2, supp\_part2 & \texttt{results/\allowbreak{}explore\_\allowbreak{}relation/\allowbreak{}snapshots/\allowbreak{}20261021-1355\_\allowbreak{}e52/\allowbreak{}e52\_\allowbreak{}belief.csv} \\
    \texttt{p2-e54-pq} & 2 & supp\_part2 & \texttt{results/\allowbreak{}explore\_\allowbreak{}relation/\allowbreak{}snapshots/\allowbreak{}20261021-1625\_\allowbreak{}e54/\allowbreak{}e54\_\allowbreak{}targets.csv} \\
    \texttt{p2-e58-lll} & 12 & supp\_part2 & \texttt{results/\allowbreak{}explore\_\allowbreak{}relation/\allowbreak{}snapshots/\allowbreak{}20261021-0440\_\allowbreak{}e49/\allowbreak{}e49\_\allowbreak{}targets.csv} \\
    \texttt{p2-e58-r} & 11 & supp\_part2 & \texttt{results/\allowbreak{}explore\_\allowbreak{}relation/\allowbreak{}snapshots/\allowbreak{}20261021-2245\_\allowbreak{}e58/\allowbreak{}e58\_\allowbreak{}targets.csv} \\
    \texttt{p2-e58-ratio} & 10 & m\_part2, supp\_part2 & \texttt{results/\allowbreak{}explore\_\allowbreak{}relation/\allowbreak{}snapshots/\allowbreak{}20261021-2245\_\allowbreak{}e58/\allowbreak{}e58\_\allowbreak{}targets.csv} \\
    \texttt{p2-e59k-a} & 5 & m\_part2, supp\_part2 & \texttt{results/\allowbreak{}explore\_\allowbreak{}relation/\allowbreak{}snapshots/\allowbreak{}20261022-0055\_\allowbreak{}e59k/\allowbreak{}e59k\_\allowbreak{}targets.csv} \\
    \texttt{p2-e59k-ctrl} & 6 & m\_part2, supp\_part2 & \texttt{results/\allowbreak{}explore\_\allowbreak{}relation/\allowbreak{}snapshots/\allowbreak{}20261022-0055\_\allowbreak{}e59k/\allowbreak{}e59k\_\allowbreak{}targets.csv} \\
    \texttt{p2-e59k-info} & 8 & supp\_part2 & \texttt{results/\allowbreak{}explore\_\allowbreak{}relation/\allowbreak{}snapshots/\allowbreak{}20261022-0055\_\allowbreak{}e59k/\allowbreak{}e59k\_\allowbreak{}targets.csv} \\
    \texttt{p2-e59k-t} & 5 & supp\_part2 & \texttt{results/\allowbreak{}explore\_\allowbreak{}relation/\allowbreak{}snapshots/\allowbreak{}20261022-0055\_\allowbreak{}e59k/\allowbreak{}e59k\_\allowbreak{}targets.csv} \\
    \bottomrule
  \end{tabularx}\par
  {\raggedright\tiny key の実体（値・書式・出所・確かめた日）は paper/results/numbers.json．数は資料で使った key の数\par}
\end{frame}
\begin{frame}{付録 4：数値の出所（10／12）}
  \rl{数値の出所}
  \vspace{-0.3em}\centering\tiny\renewcommand{\arraystretch}{1.0}
  \begin{tabularx}{\textwidth}{@{}l r l X@{}}
    \toprule key の頭 & 数 & 使った part & 台帳の出所（結果ファイル） \\ \midrule
    \texttt{p2-e59k-v} & 5 & supp\_part2 & \texttt{results/\allowbreak{}explore\_\allowbreak{}relation/\allowbreak{}snapshots/\allowbreak{}20261022-0055\_\allowbreak{}e59k/\allowbreak{}e59k\_\allowbreak{}targets.csv} \\
    \texttt{p2-e60-at} & 15 & supp\_part2 & \texttt{results/\allowbreak{}explore\_\allowbreak{}relation/\allowbreak{}snapshots/\allowbreak{}20261022-0810\_\allowbreak{}e60/\allowbreak{}e60\_\allowbreak{}targets.csv} \\
    \texttt{p2-e60-av} & 15 & m\_part2, supp\_part2 & \texttt{results/\allowbreak{}explore\_\allowbreak{}relation/\allowbreak{}snapshots/\allowbreak{}20261022-0810\_\allowbreak{}e60/\allowbreak{}e60\_\allowbreak{}targets.csv} \\
    \texttt{p2-e60-avt} & 20 & supp\_part2 & \texttt{results/\allowbreak{}explore\_\allowbreak{}relation/\allowbreak{}snapshots/\allowbreak{}20261022-0810\_\allowbreak{}e60/\allowbreak{}e60\_\allowbreak{}targets.csv} \\
    \texttt{p2-e60-ctrl} & 2 & supp\_part2 & \texttt{results/\allowbreak{}explore\_\allowbreak{}relation/\allowbreak{}snapshots/\allowbreak{}20261022-0810\_\allowbreak{}e60/\allowbreak{}e60\_\allowbreak{}targets.csv} \\
    \texttt{p2-e60-f} & 2 & m\_part2 & \texttt{results/\allowbreak{}explore\_\allowbreak{}relation/\allowbreak{}snapshots/\allowbreak{}20261022-0810\_\allowbreak{}e60/\allowbreak{}e60\_\allowbreak{}targets.csv} \\
    \texttt{p2-e60-m1} & 3 & m\_part2, supp\_part2 & \texttt{results/\allowbreak{}explore\_\allowbreak{}relation/\allowbreak{}snapshots/\allowbreak{}20261022-0810\_\allowbreak{}e60/\allowbreak{}e60\_\allowbreak{}targets.csv} \\
    \texttt{p2-e60-m2} & 1 & supp\_part2 & \texttt{results/\allowbreak{}explore\_\allowbreak{}relation/\allowbreak{}snapshots/\allowbreak{}20261022-0810\_\allowbreak{}e60/\allowbreak{}e60\_\allowbreak{}targets.csv} \\
    \texttt{p2-e60-m3} & 3 & m\_part2, supp\_part2 & \texttt{results/\allowbreak{}explore\_\allowbreak{}relation/\allowbreak{}snapshots/\allowbreak{}20261022-0810\_\allowbreak{}e60/\allowbreak{}e60\_\allowbreak{}targets.csv} \\
    \texttt{p2-e60-vt} & 15 & supp\_part2 & \texttt{results/\allowbreak{}explore\_\allowbreak{}relation/\allowbreak{}snapshots/\allowbreak{}20261022-0810\_\allowbreak{}e60/\allowbreak{}e60\_\allowbreak{}targets.csv} \\
    \bottomrule
  \end{tabularx}\par
  {\raggedright\tiny key の実体（値・書式・出所・確かめた日）は paper/results/numbers.json．数は資料で使った key の数\par}
\end{frame}
\begin{frame}{付録 4：数値の出所（11／12）}
  \rl{数値の出所}
  \vspace{-0.3em}\centering\tiny\renewcommand{\arraystretch}{1.0}
  \begin{tabularx}{\textwidth}{@{}l r l X@{}}
    \toprule key の頭 & 数 & 使った part & 台帳の出所（結果ファイル） \\ \midrule
    \texttt{p2-e60pre-juar} & 17 & appendix, m\_end, m\_part2, supp\_part2 & \texttt{results/\allowbreak{}explore\_\allowbreak{}relation/\allowbreak{}e60\_\allowbreak{}mm\_\allowbreak{}check.csv} \\
    \texttt{p2-e60pre-uar} & 4 & supp\_part2 & \texttt{results/\allowbreak{}explore\_\allowbreak{}relation/\allowbreak{}e60\_\allowbreak{}mm\_\allowbreak{}check.csv} \\
    \texttt{st-a} & 12 & m\_compare, supp\_compare & \texttt{results/\allowbreak{}stats\_\allowbreak{}compare/\allowbreak{}tests.csv} \\
    \texttt{st-b} & 3 & m\_end, supp\_compare & \texttt{results/\allowbreak{}stats\_\allowbreak{}compare/\allowbreak{}tests.csv} \\
    \texttt{st-ft} & 2 & supp\_compare & \texttt{results/\allowbreak{}stats\_\allowbreak{}compare/\allowbreak{}tests.csv} \\
    \texttt{st-oursx} & 2 & supp\_compare & \texttt{results/\allowbreak{}stats\_\allowbreak{}compare/\allowbreak{}tests.csv} \\
    \texttt{tab-at} & 2 & m\_compare & \texttt{results/\allowbreak{}baselines\_\allowbreak{}fuse-audio-text/\allowbreak{}<ターゲット>.csv} \\
    \texttt{tab-av} & 2 & m\_compare & \texttt{results/\allowbreak{}baselines\_\allowbreak{}fuse-audio-video/\allowbreak{}<ターゲット>.csv} \\
    \texttt{tab-avt} & 8 & m\_compare & \texttt{results/\allowbreak{}baselines\_\allowbreak{}fuse-audio-video-text/\allowbreak{}<ターゲット>.csv} \\
    \texttt{tab-vt} & 2 & m\_compare & \texttt{results/\allowbreak{}baselines\_\allowbreak{}fuse-video-text/\allowbreak{}<ターゲット>.csv} \\
    \bottomrule
  \end{tabularx}\par
  {\raggedright\tiny key の実体（値・書式・出所・確かめた日）は paper/results/numbers.json．数は資料で使った key の数\par}
\end{frame}
\begin{frame}{付録 4：数値の出所（12／12）}
  \rl{数値の出所}
  \vspace{-0.3em}\centering\tiny\renewcommand{\arraystretch}{1.0}
  \begin{tabularx}{\textwidth}{@{}l r l X@{}}
    \toprule key の頭 & 数 & 使った part & 台帳の出所（結果ファイル） \\ \midrule
    \texttt{tab-w2} & 8 & m\_compare & \texttt{results/\allowbreak{}baselines\_\allowbreak{}fuse2-audio-video-text/\allowbreak{}<ターゲット>.csv} \\
    \bottomrule
  \end{tabularx}\par
  {\raggedright\tiny key の実体（値・書式・出所・確かめた日）は paper/results/numbers.json．数は資料で使った key の数\par}
\end{frame}
}{}

\nocite{bendavid2010theory,schuller2010crosscorpus,busso2008iemocap,zadeh2018mosei,poria2019meld,dhall2012afew,livingstone2018ravdess,cao2014cremad,chen2022wavlm,hsu2021hubert,liu2019roberta,he2023debertav3,ryumina2022facial,radford2021clip,yolov8face,cao2018vggface2,mollahosseini2019affectnet,kumar2022lpft,gulrajani2021domainbed,okamoto1963asymptotic,raudys1998intrinsic,dobriban2018highdimensional,desilva2023value,pregibon1981logistic,koh2019accuracy,kimeldorf1971results,scholkopf2001generalized,chaloner1984optimal,chaloner1995bayesian,vapnik2009new,lopezpaz2016unifying,lee2013pseudolabel,grandvalet2004semisupervised,hinton2015distilling,blum1998combining,sindhwani2005coregularization,saerens2002adjusting,veilleux2021realistic}
\begin{frame}[allowframebreaks]{付録 5：参考文献（この資料で名前を出したもの）}
  \rl{参考文献}\label{a:refs}
  \tiny
  \bibliography{references_ieee}
\end{frame}
